\section{Introduction}

Multi-agent LLM pipelines often give their agents distinct tasks (e.g.\ critic, verifier, or
devil's advocate), and attribute observed downstream gains to the corresponding information
processing event. However, the information processing event is often not directly validated as
having occurred: the downstream metric alone is not sufficient to validate its existence, since
downstream results are a function of all aspects of the prompt, including the defined
information processing event. Prior work shows that LLMs better identify errors in others' text
as compared to their own \citep{chen2026selfcorrection,stechly2024selfverification}, but that
only justifies adding a second agent, not a particular persona.

In the experimental sciences, it is common to validate whether or not an intervention occurred
via \emph{manipulation checks}. If one cannot validate that their intervention occurred, then
there is no reason to believe they have measured the intervention. We introduce the manipulation
check to the evaluation of multi-agent LLM pipelines. Our goal is only to transfer the tool to
LLM research, not to assert its primacy: \S\ref{sec:related} names the closest existing work,
and Appendix~\ref{app:novelty-search} states what our search covered and what it did not.

In our work we observe this channel directly: an empty role label prompt refuses at a
statistically significant lower rate versus either the entire task or the length-controlled
version, with no task performed. If a model behaves differently when asked different questions,
there will be a change in score independent of whether the task was performed; in this case, the
difference lies on un-scored items, but we've still seen the mechanism in action (if not its
downstream consequences).

For this, we use six functional roles (each with a unique information processing operation)
drawn from a proposed council design. While the roles are motivated by a computational theory of
cognition \citep{lee2025emergence}, they are simply the \emph{subject of our analysis, not our
point} --- we do not endorse the council design. We find that role prompts mostly do not do what
they promise: one in six roles passes its manipulation check, one role changes the score
significantly in the \emph{opposite} direction of their stated function, and three are not
distinguishable from a length-matched control with no prescribed operation, with one
non-diagnostic after audit.

\paragraph{Contributions.}
\textbf{(i) Manipulation Checks for Role Prompting Agents}: an implementable and pre-registered
role manipulation check (compared against a length-matched neutral control), for six roles, five
of which were pre-registered with an explicit anti-degeneracy term. Post-collection stress tests
show that the presence of such a term does not establish its effectiveness: three didn't vary on
the observed support, and two admit simple degenerate solutions.
\textbf{(ii) An Analysis of Manipulation-Check-Passing Rates vs.\ Score-Changing Rates}: six
roles were categorized along both dimensions, by a pre-registered rule (the latter dimension is
null). Again, the \emph{point} is the check and criteria; the results are irrelevant.
\textbf{(iii) A Multi-Family Model Analysis}: our manipulation checks analyzed against the same
items in an entirely separate pretraining family --- shows that while our overall conclusions
remain the same, \emph{one validated check does not transfer across the two families tested},
motivating our choice to perform the manipulation check, and \emph{not} simply use a checklist.
\textbf{(iv) A Pre-Registered, Auditable Experiment}: pre-registered SESOIs; pre-registered
per-role predicted harms (and benefits); structural safeguards ensuring no fabricated data makes
it into a results table; a first-class table of exclusions.\footnote{Code, prompts, collected
records, both benchmark corpora and the pre-registration are available at
\url{https://github.com/kevzho/role-prompted-systems}. \texttt{REPRODUCE.md} gives
the exact command for each reported result; all of them run offline against the shipped records,
with no API key and no network access.}
